% =============================================================================
%  
%  
% =============================================================================

\section{Resumen Tesis}
\begin{frame}{En pocas Palabras}
  \footnotesize
  \begin{itemize}
    \setlength{\itemsep}{0pt}
    \item \textbf{Premisa:} Se evalúa el impacto de reconfigurar la red de transporte público de la \zmvm\ mediante la incorporación de un corredor anillar de alta capacidad sobre el \termino{Anillo Periférico Interior}.
    \item \textbf{Herramientas:} Se crean 3 capas de Software: Apimetro como Base de datos espaciales, VFTModel como motor analítico y Transport-Gis-Mjg para la creación de mapas
    \item \textbf{Salientes}: Para diagnosticar la red sin caer en el sesgo de compensabilidad de las métricas propuestas se propone la \termino{Clasificación Garibelt} y el \termino{Espectro MJG} donde se evalúan cinco dimensiones metodológicas.
  \end{itemize}

  \begin{block}{\textquestiondown{}API?}
    \footnotesize
    \raggedright
    El término \textbf{API} tuvo que ser creado para la división en dos secciones; para este estudio se considera al Anillo Periférico real como un Anillo Periférico Interior debido a la expansión de la \zmvm.
  \end{block}
\end{frame}
